\section{Introduction}

% several technics of PNT without GNSS

Techniques for positioning, navigation, and timing are obviously crucial not only for our personal daily life but also for the imaging of huge materials (火山の透視), automated industry, and so on. The PNT using the radio wave from the satellite, {\it e.g.} GNSS is widely used technique, however, the performance is strongly degraded if the radio waves are blocked by the obstacles such as the buildings, soils, and sea water. To overcome this problems, indoor PNT techniques are developed using the electric waves from the references such as WiFi station, laser, ... But there are different kind of obstacles may exist indoor environments or the electric waves cannot penetrate the soils nor the sea water. The application of cosmic-ray muon instead of the electric waves is performed to realize the imaging of large  objects, timing synchronization, and navigation. The muonic position system (muPS) and Muometric Wireless Navigation System  (MuWNS) was proposed by Takana based on the several time-of-flight measurements between the clients and some reference detectors assuming the velocity of the muon is the same as the speed of lights. In this method due to the clock drifting the positioning less than a few meter accuracy cannot be achieved since the frequency of the muon coincidence measurement is limited. Tanaka proposed a new concept of Vector muPS, where the angle of muon is measured by the reference and the position of the receiver is estimated by combining the angle and position measured by the reference detector. But Vector muPS does not provides the timing synchronization. 

Here we propose the new concept using the air shower where the cosmic-ray muon belongs to, toward the simultaneous positioning and timing synchronization using the cosmic-ray muon. The air showers are constantly pouring down on the ground.
The time evolution of the air showers can be estimated by measuring the shower front which travels with the speed of light. Assuming the muon is travels along the normal direction of the shower front, the virtual reference detector can be positioned at any place inside the shower front. The shower front is known to have a certain thickness for example, 100 ns at xxx m front the shower center with xxx-eV primary cosmic ray \todoi{大気シャワーの厚さについての文献を入れる}

If the detection frequency is much larger than the muon \todoi{大気シャワーフロント情報同定のエネルギー依存性の導出}, the uncertainly in the PNT can be compensated according to the statistics. For example, the frequency of coincidence with shower and muon is 10000 per second, the resolution becomes order of ten centimeter. In this paper, the efficiency of the shower front measurement and the performance of positioning and time synchronization with air-shower virtual reference are described.

